% !TEX root = chapter-standalone.tex

\section[DP as a category]{Design problems form a category}
\label{sec:dp-category}

In \cref{sec:dp-series-composition} we saw how to connect two \SY{design problems} in series.
Given $\adpa \colon \funposA \profto \resposB$ and $\adpb \colon \funposB \profto \resposC$, their series composition $\adpa \dpthen \adpb \colon \funposA \profto \resposC$ (\cref{def:dp-series-composition}) is
\begin{equation}
    (\adpa \dpthen \adpb)(\funposAopel, \resposCel)
    =
    \bigvee_{\posBel \setin \posBset}
    \adpa(\funposAopel,\resposBel)
    \booland
    \adpb(\funposBopel,\resposCel),
\end{equation}
and for every poset $\posA$ there is an identity design problem $\dpidat{\posA}$, given by $\dpidat{\posA}(\funposAel_\F{1}^\F{*}, \resposAel_\R{2}) = (\funposAel_\F{1} \posleqof{\posA} \resposAel_\R{2})$ (\cref{def:dp-identity-series}).
We also showed that series composition is associative and that the identity design problems are neutral for it (\cref{lem:dp-series-laws}).
In the diagrammatic notation, associativity says that a chain of three design problems can be read without parentheses:
\equationsag{50_assoc_1_2_3}{fig:compositionassociativity}

These are precisely the data and the axioms of a category.

\begin{definition}[Category of design problems \DP]
    \SYNDEF{category of design problems}
    \label{def:DP}
    The category of design problems \DP consists of the following constituents:
    \begin{enumerate}
        \item \emph{Objects}: The objects of $\DP$ are \SY{posets}.
        \item \emph{Morphisms}: The morphisms of $\DP$ are \SY{design problems} (\cref{def:design-problem}).
        \item \emph{Identity morphism}: The \SY{identity morphism} $\dpid_\posgenA \colon \funposA \profto \resposA$ is the identity design problem of \cref{def:dp-identity-series}.
        \item \emph{Composition operation}: Given morphisms $\adpa \colon \funposA \profto \resposB$ and $\adpb \colon \funposB \profto \resposC$, their composition $\adpa \dpthen \adpb\colon \funposA\profto \resposC$ is their series composition (\cref{def:dp-series-composition}).
    \end{enumerate}
\end{definition}

By \cref{lem:dp-series-composition}, the composite of two design problems is again a design problem, and by \cref{lem:dp-series-laws}, composition is \SY{associative} and \SY{unital}.
Therefore, $\DP$ is a category.
In \cite{fong2019}, it is called $\feas$, or $\Prof_{\Bool}$.

\begin{remark}[Design problems and feasibility relations]
    By \cref{ex:DPsAsFeasibilityRelations}, a design problem $\adpa$ carries the same information as its feasible set $\feasibleset{\adpa}$, which is a feasibility relation, that is, a monotone relation.
    By \cref{lem:dp-series-composition}, this correspondence turns series composition of design problems into composition of relations, $\feasibleset{\adpa \dpthen \adpb} = \feasibleset{\adpa} \mthen \feasibleset{\adpb}$, and it sends identity design problems to the preorder relations $\posAleq$.
    So we could equally well have described $\DP$ as the category whose objects are posets and whose morphisms are monotone relations, composed as relations.
    In the language of \cref{chap:functors}, the two descriptions give isomorphic categories.
\end{remark}

\todotextjira{359}{\bernina: @Andrea: Implement material explaining and illustrating the perspective on composing \SY{design problems} in terms of matrices and matrix multiplication.
}

\begin{gradedexercise}[\exname{ComposingDesignProblems}]
    \label{ex:ComposingDesignProblems}

    Consider the following posets, given in terms of Hasse diagrams:
    \begin{itemize}
        \item $\speed$:
              \begin{center}
                  \includesag{exam_pos_speed}
              \end{center}
        \item $\size$:
              \begin{center}
                  \includesag{exam_pos_size}
              \end{center}
        \item $\money$:
              \begin{center}
                  \includesag{exam_pos_money}
              \end{center}

        \item $\ttime$:
              \begin{center}
                  \includesag{exam_pos_time}
              \end{center}
    \end{itemize}
    Furthermore let the poset $\speed \Ptimes \size$ be equipped with the standard product partial ordering.

    Consider the design problem given by the monotone function
    \begin{equation}
        \mora\colon (\F{\speed \times \size})\op \Ptimes \R{\money} \mto \Bool
    \end{equation}
    \begin{center}
        \includesag{exam_dp_veh}
    \end{center}
    with
    \todographics{Extend makeset for multiline}
    \begin{align*}
        \mora^{-1}(\true) = & \left\{\tup{\tup{\F{\slow}, \F{\ssmall}}, \R{\cheap}}, \tup{\tup{\F{\slow}, \F{\ssmall}}, \R{\midprice}}, \tup{\tup{\F{\slow}, \F{\ssmall}}, \R{\expensive}},\right.
        \\
                            & \tup{\tup{\F{\slow}, \F{\llarge}}, \R{\midprice}}, \tup{\tup{\F{\slow}, \F{\llarge}}, \R{\expensive}} \\
                            & \left. \tup{\tup{\F{\fast}, \F{\ssmall}}, \R{\expensive}} \right\}
    \end{align*}
    along with the design problem given by the monotone function
    \begin{equation}
        \morb\colon \F{\money}\op \Ptimes \R{\ttime} \mto \Bool
    \end{equation}
    \begin{center}
        \includesag{exam_dp_saving}
    \end{center}
    with
    \begin{equation*}
        \mapb^{-1}(\true) = \makeset{\tup{\F{\cheap}, \R{\sshort}}, \tup{\F{\cheap}, \R{\llong}}, \tup{\F{\midprice}, \R{\llong}}}.
    \end{equation*}

    \begin{enumerate}
        \item Compute the series composition $\mapa \mthenof\DP \mapb$ in the category of design problems.
              The result should be a design problem described in terms of a monotone function $(\F{\speed} \Ptimes \F{\size})\op \Ptimes \R{\ttime} \mto \Bool$.
        \item We interpret elements of $\speed \Ptimes \size$ as properties of cars that Alice considers buying.
              $\money$ represents the amounts of money that she would need to buy one of said cars, and $\ttime$ represents the amounts of time Alice could spend saving money.
              The feasibility relations $\mapa$ and $\mapb$ describe what is possible for Alice.
              According to $\mapa \mthenof\DP \mapb$, is it feasible for Alice to buy a fast small car?
              If yes, will she have to work, at minimum, a long or a short amount of time in order to save enough money to buy it?
    \end{enumerate}
\end{gradedexercise}

\solutionof{ComposingDesignProblems}
\begin{gradedexercise}[\exname{DPComposition}]
    \label{ex:DPComposition}

    Consider the following posets, given in terms of Hasse diagrams:
    \begin{itemize}
        \item $\ttime$:
              \begin{center}
                  \includesag{exam_pos_time}
              \end{center}

        \item $\money$:
              \begin{center}
                  \includesag{exam_pos_money_bis}
              \end{center}

        \item $\textbf{Tech complexity}$:
              \begin{center}
                  \includesag{exam_pos_tech}
              \end{center}

        \item $\textbf{Task complexity}$:
              \begin{center}
                  \includesag{exam_pos_task}
              \end{center}
    \end{itemize}
    Let the poset $\ttime \Ptimes \money$ be equipped with the standard product partial ordering.

    You are in charge of an engineering team that should develop a robotic system.
    It is not yet clear what strategy you wish follow for realizing the project, and you wish to do a preliminary feasibility study.
    You will potentially need to make some tradeoffs between the \emph{time} that your team works on the project and the \emph{money} that you invest.
    Furthermore, you will need to decide between implementing either a cutting edge \emph{fancy} robotic system, or a \emph{simple} one.
    This will have an impact on the \emph{complexity of the tasks} that the system will be able to handle.

    Design problems capture feasibility relations.
    In your case, the feasibility relation between \textbf{Task complexity} and \textbf{Tech complexity} is given by the design problem
    \begin{equation}
        \mora\colon \F{\textbf{Task complexity}}\op \Ptimes \R{\textbf{Tech complexity}} \mto \Bool
    \end{equation}
    \begin{center}
        \includesag{exam_dp_tasktech}
    \end{center}
    with
    \begin{equation*}
        \mora^{-1}(\true) = \makeset{ \ \tup{\F{\text{high} }, \R{ \text{fancy} }}, \tup{\F{\text{medium} }, \R{ \text{fancy} }}, \tup{\F{\text{low} }, \R{ \text{fancy} }}, \tup{\F{\text{low} }, \R{ \text{simple} }}}
    \end{equation*}
    and the feasibility relation between \textbf{Tech complexity} and $\ttime \Ptimes \money$ is given by the design problem
    \begin{equation}
        \morb \colon \F{\textbf{Tech complexity}}\op \Ptimes (\R{\ttime \times \money}) \mto \Bool
    \end{equation}
    \begin{center}
        \includesag{exam_dp_tasktechmoney}
    \end{center}
    with
    \begin{align*}
        \morb^{-1}(\true) = & \{ \tup{\F{\text{fancy} }, \R{ \tup{\text{long}, \text{100K}} }}, \tup{\F{\text{simple} }, \R{ \tup{\text{long}, \text{100K}} }}, \tup{\F{\text{simple} }, \R{ \tup{\text{long}, \text{20K}} }}, \\
                            & \tup{\F{\text{simple} }, \R{ \tup{\text{short}, \text{100K}} }} \}.
    \end{align*}

    \

    Your tasks in this exercise:
    \begin{enumerate}
        \item Compute the series composition $\mapa \mthenof\DP \mapb$ in the category of design problems.
        \item Based on the previous calculation, would it be feasible to build a robotic system that can handle medium-complexity tasks if your team works only a short time on the project but invests 100K?
    \end{enumerate}
\end{gradedexercise}

\solutionof{DPComposition}
